% Apartado 2.5 - Comparacion y propuesta para el sistema embebido - BLOQUE 1
% Los numeros salen de resultados/metricas.csv, que escribe src/comparacion.py

\section{Comparación y propuesta para el sistema embebido}
\label{sec:comparacion}

\subsection{Cómo se midió}

Las tres señales que se comparan contra la referencia limpia son la contaminada
sin filtrar, la salida del filtro temporal del apartado \ref{sec:temporal} y la
salida del método espectral del apartado \ref{sec:espectral}. Las definiciones
son las del enunciado,
\begin{align}
\mathrm{MSE} &= \frac{1}{L}\sum_{n=0}^{L-1}\bigl(s[n]-\hat{s}[n]\bigr)^2,\\
\mathrm{SNR}_{\mathrm{dB}} &= 10\log_{10}
  \left(\frac{\sum_{n} s^2[n]}{\sum_{n}(s[n]-\hat{s}[n])^2}\right),\\
\Delta\mathrm{SNR} &= \mathrm{SNR}_{\mathrm{salida}}-\mathrm{SNR}_{\mathrm{entrada}}.
\end{align}

Las tres las calcula \texttt{src/comparacion.py} leyendo los archivos de
\texttt{resultados/salidas/}, no llamando al código de cada método. Así los dos
se miden con la misma vara y ninguno puede reportarse a sí mismo.

Antes de medir se estima por correlación cruzada el retardo entre cada salida y
la referencia, y se compensa. Sin esa compensación el MSE castigaría el retardo
de grupo del filtro en vez de su distorsión. En las seis salidas entregadas el
retardo detectado fue de cero muestras, es decir que ambos métodos ya devuelven
la señal alineada.

La columna de energía útil conservada es la fracción de la energía que sobrevive
en una ventana de $\pm 60$\,Hz alrededor de cada tono. La de atenuación mide la
caída de energía en la banda donde vive el ruido de cada caso: el tono $\pm
50$\,Hz en el primero, la banda 800--1500\,Hz en el tercero, y todo lo que queda
por encima de la banda útil en el segundo, que por ser ruido blanco no tiene
banda propia.

\subsection{Resultados}

\begin{table}[H]
\centering\small
\caption{Comparación de los tres casos. Generada por \texttt{src/comparacion.py}
en \texttt{resultados/metricas.csv}.}
\begin{tabular}{llrrrrr}
\toprule
Caso & Método & MSE & SNR (dB) & $\Delta$SNR (dB) & E. útil & Atenuación (dB) \\
\midrule
\texttt{caso\_1} & sin filtrar & $4{,}88\cdot10^{-2}$ & 2,00 & 0,00 & 1,000 & 0,00 \\
\texttt{caso\_1} & temporal    & $8{,}75\cdot10^{-5}$ & \textbf{29,46} & \textbf{+27,46} & 1,000 & \textbf{65,22} \\
\texttt{caso\_1} & espectral   & $2{,}97\cdot10^{-4}$ & 24,15 & +22,15 & 1,000 & 25,33 \\
\midrule
\texttt{caso\_2} & sin filtrar & $1{,}89\cdot10^{-2}$ & $-0{,}00$ & 0,00 & 1,051 & 0,00 \\
\texttt{caso\_2} & temporal    & $1{,}82\cdot10^{-2}$ & 0,18 & +0,18 & 1,050 & 0,23 \\
\texttt{caso\_2} & espectral   & $5{,}34\cdot10^{-3}$ & \textbf{5,50} & \textbf{+5,50} & \textbf{0,867} & 6,11 \\
\midrule
\texttt{caso\_3} & sin filtrar & $1{,}94\cdot10^{-2}$ & 3,00 & 0,00 & 1,159 & 0,00 \\
\texttt{caso\_3} & temporal    & $1{,}94\cdot10^{-2}$ & 3,00 & +0,00 & 1,158 & 0,00 \\
\texttt{caso\_3} & espectral   & $1{,}73\cdot10^{-2}$ & 3,49 & +0,49 & 1,160 & 0,47 \\
\bottomrule
\end{tabular}
\end{table}

El tiempo de proceso del método espectral, medido sobre los 63 bloques de una
señal completa, fue de 3,73\,ms por bloque en una computadora de escritorio con
la FFT propia escrita en Python. El presupuesto en tiempo real es de 32\,ms por
bloque, que es lo que dura el salto de 512 muestras a 16\,kHz. El filtro temporal
cuesta 15 multiplicaciones-acumulaciones por muestra, es decir 240\,000\,MAC/s a
esa frecuencia, según la medición del apartado \ref{sec:temporal}. La medición
sobre el microcontrolador corresponde a la siguiente etapa del proyecto.

\subsection{Cuál ofrece el mejor compromiso}

No hay un ganador único, y esa es la conclusión más útil de la tabla.

En el primer caso gana el filtro temporal con holgura: 27,5\,dB de ganancia
contra 22,2\,dB, conservando ambos el 100\,\% de la energía útil. La diferencia
grande está en la atenuación de la interferencia, 65,2\,dB contra 25,3\,dB. La
razón es estructural: una muesca IIR puede hacerse arbitrariamente profunda sobre
una frecuencia fija, mientras que la atenuación por bines tiene como piso la fuga
de la ventana Hann, que reparte la energía del tono entre varios bines vecinos.
Contra una interferencia tonal estacionaria y conocida, el filtro temporal es
mejor herramienta.

En el segundo caso se invierte por completo: el método espectral da 5,5\,dB y el
temporal 0,18\,dB. Pero el espectral lo paga perdiendo el 13,3\,\% de la energía
útil, el único caso de la tabla donde esa columna baja de 1,000.

En el tercero ninguno sirve. El temporal no mueve la aguja y el espectral gana
0,49\,dB, que es menos que la incertidumbre con la que alguien juzgaría la señal
de oído.

Si hay que escoger uno solo, el espectral ofrece el mejor compromiso
\emph{promedio}, porque es el único que hace algo en los tres escenarios: su
clasificación por bloque le permite cambiar de estrategia, mientras que el
temporal tiene una $H(z)$ fija diseñada para un caso. Pero el sistema distribuido
del proyecto no obliga a escoger: tiene un nodo para cada uno, y la tabla
justifica esa decisión de arquitectura en lugar de contradecirla.

\subsection{En qué condición falla cada método}

\textbf{El filtro temporal falla cuando el ruido entra en su banda de paso.} Un
sistema lineal e invariante en el tiempo solo puede discriminar por frecuencia:
su respuesta $H(e^{j\omega})$ multiplica cada componente por un número que
depende únicamente de $\omega$. Como la banda de paso tiene que contener los
tonos útiles, cualquier ruido que caiga ahí se multiplica por el mismo factor que
la señal y sale intacto. Eso es exactamente lo que ocurre en el segundo caso,
donde el ruido blanco cubre la banda útil, y en el tercero, donde la banda de
ruido la contiene. Sus ganancias de 0,18 y 0,00\,dB no son un defecto de diseño
sino el límite de la clase de sistemas a la que pertenece.

\textbf{El método espectral falla por dos razones distintas.} La primera es la
misma del caso anterior llevada al dominio de la frecuencia: cuando señal y ruido
comparten bin, atenuar el bin atenúa las dos cosas, y de ahí los 0,49\,dB del
tercer caso y la caída de energía útil a 0,867 en el segundo. La segunda es
propia del método: la profundidad de la atenuación está limitada por la fuga de
la ventana. Por eso, incluso en el caso que le resulta fácil, llega a 25\,dB de
atenuación donde el filtro temporal llega a 65.

En resumen, el temporal falla por rigidez y el espectral por resolución. Las dos
fallas se vuelven la misma cuando señal y ruido ocupan la misma frecuencia, que
es el límite que el apartado \ref{sec:limite} ya había anticipado midiendo la SNR
tono por tono antes de filtrar nada.

\subsection{Tamaño de bloque, muestreo y formato numérico}

\begin{center}\small
\begin{tabular}{lll}
\toprule
Parámetro & Propuesta & Justificación \\
\midrule
Frecuencia de muestreo & 16\,kHz & la componente más alta del banco es de 3500\,Hz \\
Tamaño de bloque & 1024 muestras & $\Delta f = 15{,}6$\,Hz, potencia de dos \\
Solapamiento & 50\,\% (salto de 512) & reconstrucción sin artefactos con Hann \\
Latencia por bloque & 64\,ms de ventana, 32\,ms de salto & \\
Formato numérico & \texttt{float32} & el estudio de cuantización lo respalda \\
\bottomrule
\end{tabular}
\end{center}

Los 16\,kHz dejan un factor de 4,6 sobre la componente más alta del banco de
pruebas, margen de sobra para el filtro antialias analógico. Bajar a 8\,kHz
dejaría la interferencia de 3500\,Hz a solo 0,6\,kHz de Nyquist y obligaría a un
antialias mucho más abrupto; subir a 44,1\,kHz triplicaría el costo sin aportar
banda útil.

El bloque de 1024 es el punto donde la resolución de 15,6\,Hz alcanza para
separar los tonos de 440 y 880\,Hz con holgura, y la latencia de 64\,ms todavía
es tolerable. Duplicarlo a 2048 daría 7,8\,Hz pero 128\,ms de retardo, que en un
sistema que procesa audio ya se percibe.

Sobre el formato, el estudio de cuantización del apartado \ref{sec:temporal}
muestra que con \texttt{float32} la respuesta del filtro queda intacta, con 14
bits fraccionarios la desviación en banda de paso es de 0,010\,dB, y con 8 bits
el rizado casi se duplica y la atenuación cae por debajo de la especificación. La
recomendación es \texttt{float32} siempre que el procesador tenga unidad de punto
flotante; si no la tuviera, el formato Q2.14 con acumulador de mayor ancho es
suficiente, y Q1.15 queda descartado porque $|a_1|$ puede superar la unidad.

\subsection{Hardware y bibliotecas disponibles}

Se propone el \textbf{ESP32-WROOM-32} para los tres nodos. Pesan tres razones: el
grupo ya tiene unidades y experiencia con él de un proyecto anterior del curso;
trae unidad de punto flotante de precisión simple, que es lo que el formato
propuesto pide; y sus 520\,kB de SRAM dejan el presupuesto de memoria del nodo
espectral, estimado más abajo en unos 22\,kB, en territorio cómodo.

\begin{table}[H]
\centering\small
\caption{Bibliotecas consideradas para la FFT y el filtrado.}
\begin{tabular}{llp{6.4cm}}
\toprule
Biblioteca & Licencia & Precisión y limitaciones \\
\midrule
ESP-DSP \cite{espdsp} & Apache 2.0 & Biblioteca oficial de Espressif para ESP32 y ESP32-S3. Trae FFT, FIR e IIR con implementaciones en coma flotante de 32 bits y en enteros de 16 bits, varias escritas en ensamblador para el núcleo Xtensa, con referencias en C ANSI. Las rutas optimizadas están atadas a ese núcleo, de modo que migrar a otra arquitectura obliga a volver a las versiones en C. \\
\addlinespace
arduinoFFT \cite{arduinofft} & GPL-3.0 & Más simple de integrar desde el entorno de Arduino, sin optimización en ensamblador. Su licencia es \emph{copyleft}, lo que condiciona la distribución del firmware que la enlace; para un trabajo académico no es obstáculo, pero sí es un criterio a considerar. \\
\addlinespace
Implementación propia & --- & \texttt{src/dft\_fft.py} ya tiene la FFT radix-2 iterativa con reordenamiento bit-reverso, verificada contra la de biblioteca en \texttt{pruebas/test\_fft.py}. Portarla a C es directo y elimina toda dependencia externa, a costa de perder las optimizaciones en ensamblador. \\
\bottomrule
\end{tabular}
\end{table}

CMSIS-DSP, que sería la opción natural en un Cortex-M, no aplica al ESP32 por ser
de arquitectura Xtensa; quedaría sobre la mesa solo si el grupo cambiara a un
STM32.

La recomendación es ESP-DSP por la licencia permisiva y las rutinas optimizadas,
conservando la implementación propia como referencia de verificación: el nodo
puede correr ambas sobre el mismo bloque y comparar, igual que el apartado
\ref{sec:temporal} compara la ecuación de diferencias propia contra la de
biblioteca.

\subsection{Arquitectura de los tres nodos}

\begin{center}\small
\begin{tabular}{lp{10.2cm}}
\toprule
Nodo & Función \\
\midrule
1. Adquisición & Muestrea a 16\,kHz, elimina la componente de continua, aplica la ventana Hann, calcula la FFT y clasifica el ruido con las tres características del apartado \ref{sec:senales}. Reparte a los nodos 2 y 3. \\
2. Filtrado temporal & Aplica la ecuación de diferencias del apartado \ref{sec:temporal} y entrega su salida. \\
3. Filtrado espectral y evaluación & Aplica la cadena del apartado \ref{sec:espectral}, calcula MSE, SNR y energía, y muestra la comparación. \\
\bottomrule
\end{tabular}
\end{center}

\textbf{Qué se transmite.} Por cada bloque, el nodo 1 envía las 512 muestras
nuevas de la señal contaminada y un mensaje corto de parámetros: la clase de
ruido detectada (1 byte), la frecuencia interferente estimada o los dos bordes de
banda (2 o 4 bytes) y el nivel de piso de ruido (4 bytes). Los nodos 2 y 3 hacen
su propio procesamiento sobre las muestras; los parámetros les ahorran repetir la
clasificación y garantizan que los dos actúen sobre el mismo diagnóstico.

\textbf{Cuántos bytes.} Con bloques de 1024 muestras, salto de 512 y 16 bits por
muestra:
\[
512\ \text{muestras}\times 2\ \text{B} = 1024\ \text{B por bloque},
\qquad
\frac{16\,000}{512} = 31{,}25\ \text{bloques/s}.
\]
El flujo de muestras por enlace es entonces
\[
1024\ \text{B}\times 31{,}25\ \text{s}^{-1} = 32\,000\ \text{B/s} = 256\ \text{kbit/s},
\]
que es exactamente $f_s$ por los 16 bits de cada muestra, como tiene que ser. El
mensaje de parámetros añade unos 16\,B por bloque, es decir 500\,B/s o 4\,kbit/s:
menos del 2\,\% del total.

\textbf{Tasa mínima.} Cada enlace necesita 260\,kbit/s de carga útil. El nodo 1
alimenta dos nodos, de modo que su salida total es de 520\,kbit/s; los nodos 2 y
3 devuelven su señal filtrada al nodo de comparación con otros 256\,kbit/s cada
uno. Sobre UART con formato 8N1 el encuadre añade un 25\,\%, lo que sube el
requisito por enlace a unos 325\,kbit/s: un UART a 460\,800 baudios alcanza con
margen y uno a 921\,600 deja holgura para reintentos. Si se prefiere un solo bus
compartido, SPI a 1\,MHz o más cubre los tres enlaces sin acercarse al límite.

\textbf{Memoria estimada del nodo espectral.} Dos arreglos de 1024 valores en
\texttt{float32} para la parte real e imaginaria son 8\,kB; la tabla de la
ventana Hann, 4\,kB; el acumulador de solapamiento, 2\,kB; y los búferes de
entrada y salida, unos 8\,kB. El total ronda los 22\,kB, holgado frente a los
520\,kB de SRAM del ESP32. El nodo temporal necesita mucho menos: 84\,B de
coeficientes y estados según el apartado \ref{sec:temporal}, más sus búferes.

\textbf{Latencia de extremo a extremo.} El bloque impone 64\,ms de ventana más
32\,ms de salto. A eso se suma el tiempo de proceso, que en el presupuesto de
32\,ms por bloque deja margen suficiente según la medición de referencia, y el de
transmisión, que a 460\,800 baudios es de unos 23\,ms por bloque de 1024\,B. El
total queda por debajo de los 150\,ms, compatible con la operación en tiempo
cuasi real que el proyecto pide.
